% !TeX root = ../main-4.tex
\section{Related work}\label{sec:related-work}

\greenrevision{We review audio--language models, sound semantic descriptions and textual prompting, and structured knowledge for audio classification. These research directions address cross-modal matching, the information supplied by class descriptions, and the use of relations in classification, respectively.}

\subsection{Zero-shot audio classification with audio--language models}

\greenrevision{Zero-shot audio classification relates audio representations to class semantics to recognize classes without training audio from the target classes. Xie et al. learn an acoustic--semantic projection through a bilinear compatibility framework, using semantic representations from pretrained language models~\cite{xie2021zero}. Cross-modal pretraining provides a complementary route to natural-language supervision. CLIP learns transferable visual representations from image--text pairs~\cite{radford2021learning}. AudioCLIP incorporates an audio encoder into CLIP, whereas Wav2CLIP learns audio representations by distilling CLIP's visual representations~\cite{guzhov2021audioclip,wu2022wav2clip}. Both support matching audio with images and text in a shared embedding space.}

\greenrevision{Direct audio--text pretraining further develops this approach. CLAP uses two encoders and contrastive learning to map audio and text descriptions into a joint embedding space~\cite{elizalde2023clap}. LAION-CLAP expands the training data and introduces feature fusion and keyword-to-caption augmentation~\cite{wu2023clap}. M2D-CLAP combines masked audio modeling with audio--text alignment, while SLAP supports variable audio durations through joint contrastive, self-supervised, and captioning objectives~\cite{niizumi2024m2dclap,mei2026slap}. These studies strengthen the representations used for cross-modal matching. The semantic content of class descriptions also influences prediction, motivating the description and prompting methods reviewed below.}

\subsection{Sound semantic description and textual prompting}

\greenrevision{Sound semantic descriptions supply information beyond class names. Binary Attribute Embeddings represents classes through attributes such as pitch, duration, and source material~\cite{lin2022binary}. Sound Attribute Vector learns discriminative global features alongside spectrotemporal local features that predict attribute information~\cite{lin2023soundattribute}. Xu et al. use the knowledge of large language models to generate detailed auditory attribute descriptions for each class~\cite{xu2024soundattribute}. ReCLAP extends natural-language descriptions by training CLAP on rewritten audio captions and augmenting class prompts with sound descriptions in diverse scenes~\cite{ghosh2024reclap}. Attribute representations and textual descriptions thus provide different ways to supplement the semantics of short labels.}

\greenrevision{Prompting methods also consider task information and cross-modal alignment. TSPE constructs task-specific prompt ensembles using labels, sound attributes, and sound sources, then averages the prompt embeddings for each class~\cite{anand2025tspe}. For each prompt, PAT sums the maximum class prediction logit across audio samples from the downstream task. It uses these sums to weight and ensemble text representations. Its cross-modal alignment module uses attention between frame-level audio representations and text~\cite{seth2025pat}. In the related domain of sound design, AudioCards structures acoustic attributes and sonic descriptors as metadata for text--audio retrieval, descriptive captioning, and metadata generation~\cite{sridhar2026audiocards}. PAT's prompt weighting uses information across a task, while AudioCards addresses sound-design tasks rather than zero-shot class reranking.}

\greenrevision{These methods enrich attributes, descriptions, or prompt representations. Graph triples introduce a further requirement: the text needs to express the relation connecting the head and tail entities. Concatenating their names leaves this relation unstated, whereas unconstrained expansion may add content absent from the triple. AAKV therefore verbalizes the supplied entities and directed relation as a description of sound under explicit content constraints.}

% Queue the method overview early so it appears at the top of page 3.
\providecommand{\methodframeworkfigure}{figures/saki_framework.png}
\providecommand{\methodframeworkwidth}{0.96\textwidth}
\begin{figure*}[!t]
  \centering
  \includegraphics[width=\methodframeworkwidth]{\methodframeworkfigure}
  \caption{Overview of SAKI. Offline, class-to-entity mapping, RotatE tail retrieval, and AAKV produce cached graph evidence and CLAP text embeddings. Online, (a) CLAP identifies the top-$K$ classes; (b) \redrevision{class rankings computed separately for each relation provide top-1 predictions and top-two score gaps for sample-level relation selection}; and (c) the selected evidence is deduplicated and aggregated by class before fusion with the original CLAP score. Classes outside the initial top-$K$ or without usable evidence retain their original CLAP scores.}
  \label{fig:framework}
\end{figure*}

\subsection{Structured knowledge for audio classification}

\greenrevision{Structured knowledge represents class hierarchies and relations among audio events. AudioSet organizes sound classes through a hierarchical ontology~\cite{gemmeke2017audioset}. Jimenez et al. design neural architectures that preserve ontology structure~\cite{jimenez2018ontology}. Sun et al. combine feed-forward ontology layers with graph convolution to model label dependencies across and within ontology levels~\cite{sun2020ontology}. Aironi et al. compare node embeddings in a classification framework whose label graph is derived from label co-occurrence in training data~\cite{aironi2022graph}. Hou et al. construct event relational graphs for acoustic scene classification, representing events as nodes and pairwise relationship cues as multidimensional edge features~\cite{hou2023audio}. These approaches incorporate class or event structure into model learning; retrieving external graph evidence during audio--language inference involves a different process for using knowledge.}

\greenrevision{Knowledge graph embeddings enable relational retrieval through triple scores. RotatE represents each relation as a rotation from the head entity to the tail entity in complex space for link prediction~\cite{sun2019rotate}. iKnow-audio applies this retrieval process to CLAP's initial class predictions using a curated subset of informative relations. It concatenates class labels with retrieved tail entities to form enriched prompts, then jointly aggregates the original and enriched prompt scores through log-sum-exp~\cite{olvera-etal-2025-iknow}. Its graph scores assess entity--relation compatibility, while its initial candidates depend on the audio. SAKI further evaluates relations using predictions for the current audio and separates aggregation of knowledge evidence from fusion with the original CLAP score.}

\greenrevision{Input-dependent knowledge selection has also been studied in image classification. SINet selects graph knowledge through an attention-based module and uses a knowledge selection loss to measure consistency between knowledge and image features~\cite{tang2023sinet}. It offers a related example of selecting knowledge for an input, but learns the selection module within an image classification framework. SAKI selects relations using top-1 voting and top-two score gaps without training a selector on the target datasets.}

\greenrevision{Together, these research directions provide cross-modal representations, richer class semantics, and explicit relational knowledge. SAKI builds on this foundation to organize graph knowledge for individual audio samples at inference. The Sample-level Relation Selector determines which relations contribute evidence, AAKV expresses their triples under content constraints, and Hierarchical Fusion combines that evidence with the original CLAP prediction.}
